\documentclass[11pt]{article}
\usepackage[margin=1in]{geometry}
\usepackage{amsmath,amssymb,amsthm}
\usepackage[hidelinks]{hyperref}
\newtheorem{theorem}{Theorem}
\newtheorem{lemma}{Lemma}
\newcommand{\R}{\mathcal R}
\title{The strict resultant comparison for the \(L_\alpha\) means}
\author{Principia Math In Progress}
\date{18 September 2026}
\begin{document}
\maketitle
\begin{abstract}
We prove the full nondegenerate range of the \(L_\alpha\)
resultant-mean sign conjecture for both power-mean branches.
The proof reduces a mixed derivative to strict concavity of
\(\sqrt u\coth(t\sqrt u)\) in \(u\). The zero-parameter case is
the already known logarithmic-mean equality, which we state explicitly.
\end{abstract}

\section*{Task 1: the original problem}
For positive \(a\ne b\), the source defines
\[
 L_\alpha(a,b)=\frac{2\alpha a^\alpha b^\alpha(b-a)}
                         {b^{2\alpha}-a^{2\alpha}},
 \qquad
 L_0(a,b)=\frac{b-a}{\log b-\log a};
 \qquad L_\alpha(a,a)=a.
\]
These means are symmetric, homogeneous, and satisfy
\(L_{-\alpha}=L_\alpha\).
Let \(B_r(a,b)=((a^r+b^r)/2)^{1/r}\) for \(r\ne0\),
and \(B_0(a,b)=\sqrt{ab}\). Write
\[
 \R(B_p,L_\alpha,B_q)(a,b)
 =B_p\bigl(L_\alpha(a,B_q(a,b)),L_\alpha(B_q(a,b),b)\bigr).
\]
Conjecture~4.2 of the preprint \cite{MR}, retained as Conjecture~32
on p.~20 of the complete published article, asks for the negative
sign of \(L_\alpha-\R\) at the following parameters.

\begin{theorem}\label{main}
Suppose \(0<|\alpha|<1/2\), and put
\[
 r=\sqrt{1+16\alpha^2-16\alpha^4},\qquad
 p=r\ \text{or}\ p=-r,\qquad
 q=\frac{1-p}{2}-2\alpha^2.
\]
Then, for every distinct positive pair \(a,b\),
\[
 L_\alpha(a,b)<\R(B_p,L_\alpha,B_q)(a,b).
\]
At \(\alpha=0\), the same parameters give equality for all positive
inputs. Equality also holds at \(|\alpha|=1/2\) with the limiting
parameter pairs, and on the diagonal \(a=b\).
\end{theorem}

The source explicitly records the logarithmic-mean equality at
\(\alpha=0\) in the example preceding its conjecture.
Thus the printed strict inequality's inclusion of zero is an
already known degenerate exception. Our contribution is the proof
of the entire nonzero range, not that exception.

\section*{A strict-concavity lemma}
\begin{lemma}\label{concavity}
For fixed \(t>0\), the function
\[
 \Phi_t(u)=\sqrt u\,\coth(t\sqrt u),\quad u>0,\qquad
 \Phi_t(0)=1/t
\]
is strictly increasing and strictly concave on \([0,\infty)\).
\end{lemma}
\begin{proof}
The classical partial-fraction expansion for \(\coth\),
\cite[Eq.~4.36.3]{DLMF}, gives
\[
 \Phi_t(u)=\frac1t+
       2t\sum_{n=1}^{\infty}\frac{u}{\pi^2n^2+t^2u}.
\]
Each summand is increasing and strictly concave. More explicitly,
\[
 \Phi_t'(u)=2t\sum_{n=1}^{\infty}
       \frac{\pi^2n^2}{(\pi^2n^2+t^2u)^2}>0,\qquad
 \Phi_t''(u)=-4t^3\sum_{n=1}^{\infty}
       \frac{\pi^2n^2}{(\pi^2n^2+t^2u)^3}<0.
\]
The differentiated series converge uniformly on compact nonnegative
intervals, being bounded by constant multiples of \(n^{-2}\) and
\(n^{-4}\), respectively. This proves the lemma, including zero by
continuity.
\end{proof}

\section*{The parameter identity}
By symmetry in \(\alpha\), assume \(0<\alpha<1/2\). Put
\[
 A=\frac12,\quad B=\alpha,\quad
 \kappa=1-4\alpha^2,\quad C=\frac{|q|}{2}.
\]
Here \(0<\kappa<1\), \(1<r<2\), and
\begin{equation}\label{circle}
 r^2+\kappa^2+2\kappa=4,\qquad q=(\kappa-p)/2.
\end{equation}
In particular \(pq<0\).
For each \(t>0\), abbreviate
\(f_c=c\coth(ct)=\Phi_t(c^2)\).

If \(p=r\), then \(C=(r-\kappa)/4>0\), and \eqref{circle} gives
\[
 2C^2=(r+1)B^2-(r-1)A^2.
\]
Thus \(B^2\) is a strict convex combination of \(C^2\) and \(A^2\):
\[
 B^2=\frac2{r+1}C^2+\frac{r-1}{r+1}A^2.
\]
The points are distinct because \(B<A\). Strict concavity gives
\begin{equation}\label{pos}
 2f_C<(r+1)f_B-(r-1)f_A.
\end{equation}

If \(p=-r\), then \(C=(r+\kappa)/4\), and instead
\[
 2C^2=(r+1)A^2-(r-1)B^2.
\]
Applying the same strict concavity with \(A^2\) as the intermediate
point gives
\begin{equation}\label{neg}
 2f_C<(r+1)f_A-(r-1)f_B.
\end{equation}
Equivalently, both cases establish the single strict inequality
\begin{equation}\label{W}
 W(t):=(p-1)f_A-(p+1)f_B+2f_C<0.
\end{equation}

\section*{From concavity to the resultant mean}
Let
\[
 v(t)=\log\frac{2\alpha\sinh(t/2)}{\sinh(\alpha t)}.
\]
Then
\[
 v'=f_A-f_B>0
\]
by the increasing part of Lemma~\ref{concavity}.
Since \(f_c'=c^2-f_c^2\), we also have
\[
 v''=\frac\kappa4-(f_A+f_B)v'.
\]
Consequently
\begin{align}
 T(t)&:=v''+p(v')^2+q\coth(qt/2)v'-\frac\kappa4
       \nonumber\\
 &=v'\bigl((p-1)f_A-(p+1)f_B+2f_C\bigr)<0.
 \label{T}
\end{align}
Here \(q\coth(qt/2)=2C\coth(Ct)=2f_C\), for either sign of \(q\).

For \(a<b\), define
\[
 H(a,b)=(b^q-a^q)L_\alpha(a,b)^p,\qquad H(a,a)=0.
\]
Writing \(s=p+q\), \(t=\log(b/a)\), we obtain
\[
 H(a,b)=a^s h(t),\qquad
 h(t)=e^{st/2}\,2\sinh(qt/2)e^{pv(t)}.
\]
Differentiation, using \(p+2q=\kappa\), gives
\[
 h''-s h'=p hT,\qquad
 H_{ab}(a,b)=-\frac{a^s}{ab}\,p h(t)T(t)<0.
\]
Indeed \(h\) has the sign of \(q\), \(pq<0\), and \(T<0\).
It follows, for every \(a<c<b\), that
\begin{equation}\label{rectangle}
 H(a,c)+H(c,b)-H(a,b)
   =\int_a^c\int_c^b H_{ab}(u,v)\,dv\,du<0.
\end{equation}
Take \(c=B_q(a,b)\), so
\(c^q-a^q=b^q-c^q=(b^q-a^q)/2\).
If \(q>0\), \eqref{rectangle} gives
\[
 L_\alpha(a,b)^p>
     \frac{L_\alpha(a,c)^p+L_\alpha(c,b)^p}{2}.
\]
Now \(p<0\), so taking the \(p\)-th root proves \(L_\alpha<\R\).
If \(q<0\), division by \(b^q-a^q\) reverses the displayed
inequality, and \(p>0\); again \(L_\alpha<\R\).
Symmetry covers \(a>b\).

\paragraph{The equality cases.}
At \(\alpha=0\), the parameter pairs are \((p,q)=(1,0)\) and
\((-1,1)\). Splitting the logarithmic mean at \(c=\sqrt{ab}\) gives
\[
 \frac{L_0(a,c)+L_0(c,b)}2=L_0(a,b).
\]
Splitting at \(c=(a+b)/2\) gives
\[
 \frac{L_0(a,c)^{-1}+L_0(c,b)^{-1}}2=L_0(a,b)^{-1}.
\]
These are the already known equalities cited by the source.
At \(|\alpha|=1/2\), \(L_\alpha=\sqrt{ab}\) and the parameter pairs
are \((2,-1)\) and \((-2,1)\). Direct substitution, using
\(B_1(a,b)B_{-1}(a,b)=ab\), gives equality.
The diagonal equality is immediate from the definitions.
This completes Theorem~\ref{main}. \(\square\)

\paragraph{Verification and scope.}
The proof above is entirely analytic. The accompanying verifier checks
six exact algebraic identities, and supplements them with 126 strict
comparisons and 36 endpoint equalities at 100 decimal digits.
Those finite checks audit the formulas and signs; they are not the
universal proof. The authoring agent conducted the audit; independent
peer review is not claimed.

\section*{Proposed follow-up questions}
The following four tasks are unproved directions, not additional solves.
\paragraph{Task 2.}
For each fixed \(0<\alpha<1/2\), determine the full region of
parameter pairs \((p,q)\) for which \(L_\alpha<\R(B_p,L_\alpha,B_q)\)
holds for every distinct positive input pair.
\paragraph{Task 3.}
For each of the two parameter branches, determine the sharp
constant \(K_\alpha>0\), if it exists, in
\[
 \frac{\R(B_p,L_\alpha,B_q)-L_\alpha}{L_\alpha}
 \ge K_\alpha
 \left(\frac{|\log(b/a)|}{1+|\log(b/a)|}\right)^6.
\]
\paragraph{Task 4.}
For \(1/2<|\alpha|<1\), determine the number and location of sign changes
of the same resultant difference as \(|\log(b/a)|\) varies.
\paragraph{Task 5.}
Characterize which homogeneous means with a hyperbolic-product formula
for their associated one-variable function admit a resultant comparison
deduced from concavity in squared frequencies.

\begin{thebibliography}{9}
\raggedright
\bibitem{MR}
L.~Mihokovi\'c and M.~Ra\"issouli,
\emph{Results on comparison and sub/super-stabilizability of some new means},
\href{https://arxiv.org/abs/2405.13618}{arXiv:2405.13618v2} (2024),
Conjecture~4.2; Applicable Analysis and Discrete Mathematics
\textbf{20} (2026), 1--23,
\href{https://doi.org/10.2298/AADM240618001M}{doi:10.2298/AADM240618001M},
Conjecture~32.
Complete published text:
\href{http://pefmath.etf.rs/vol20num1/AADM-Vol20-No1-1-23.pdf}{journal PDF}.
\bibitem{DLMF}
NIST Digital Library of Mathematical Functions,
\href{https://dlmf.nist.gov/4.36.E3}{Eq.~4.36.3},
the classical partial-fraction expansion of the hyperbolic cotangent.
\end{thebibliography}
\end{document}
